\section{Análise Experimental}

A avaliação experimental da plataforma Wumpus Verse foi conduzida com o propósito de verificar o correto funcionamento de suas principais funcionalidades e a integração entre os módulos que compõem o sistema. Para isso, foram executados testes funcionais baseados em cenários representativos das ações que um usuário realiza durante a utilização da plataforma, abrangendo desde o cadastro e a autenticação até a criação de mundos e agentes, a execução de simulações, o armazenamento de resultados e as operações de importação e exportação de cenários. Dessa forma, buscou‑se assegurar que o comportamento observado estivesse em conformidade com os requisitos definidos para o protótipo.

\subsection{Objetivo da Validação}

A validação teve como objetivo principal verificar o funcionamento das funcionalidades disponibilizadas pela plataforma, avaliando sua capacidade de executar corretamente as operações previstas e de integrar os diferentes módulos do sistema. Além disso, procurou‑se identificar eventuais limitações ou aspectos que ainda necessitassem de ajustes, garantindo que os usuários pudessem realizar as tarefas essenciais do fluxo de utilização sem a ocorrência de falhas críticas. Para tanto, foram considerados aspectos relacionados à autenticação, ao gerenciamento de mundos, à criação e configuração de agentes, à execução e ao armazenamento de simulações, à geração de estatísticas e à manipulação de dados por meio dos recursos de importação e exportação.

\subsection{Metodologia de Testes}

A validação da plataforma foi realizada por meio de testes funcionais manuais baseados em cenários de uso. Essa abordagem foi adotada por ser adequada para avaliar o comportamento das funcionalidades implementadas em uma aplicação que ainda se encontra em fase de prototipação, permitindo verificar diretamente os resultados obtidos a partir das ações executadas pelo usuário. Para cada cenário de teste, foi definida uma funcionalidade alvo, um critério de sucesso correspondente ao comportamento esperado e o resultado observado durante a execução. Os cenários foram elaborados de modo a contemplar as principais etapas de utilização da plataforma, abrangendo o gerenciamento de usuários, ambientes e agentes, bem como a execução, a consulta e a manipulação dos resultados das simulações.

A Tabela 2 apresenta os cenários utilizados durante a validação, juntamente com os respectivos critérios de sucesso e resultados observados.

\begin{table}[ht]
    \centering
    \caption{Apresenta os cenários de teste.}
    \label{tab:cenario-testes}
    \small
    \begin{tabular}{|p{1cm}|p{4.5cm}|p{6.5cm}|p{2cm}|}
        \hline
        \textbf{ID} & \textbf{Funcionalidade} & \textbf{Critério de Sucesso}  & \textbf{Resultado} \\
        \hline
        CT01 & Cadastro de usuário & Usuário cadastrado com sucesso & Funcional \\
        \hline
        CT02 & Login na plataforma & Autenticação realizada corretamente & Funcional \\
        \hline
        CT03 & Criação de mundo & Mundo criado e armazenado sem erros & Funcional \\
        \hline
        CT04 & Criação de agentes & Agente criado e salvo corretamente & Funcional \\
        \hline
        CT05 & Funcionamento dos agentes & Agente realiza ações coerentes no ambiente & Ajustes \hbox{Planejados} \\
        \hline
        CT06 & Execução de simulação & Simulação executada e resultados apresentados ao usuário & Ajustes \hbox{Planejados} \\
        \hline
        CT07 & Visualização do histórico & Execuções exibidas corretamente & Funcional \\
        \hline
        CT08 & Geração de estatísticas & Métricas apresentadas corretamente & Ajustes \hbox{Planejados} \\
        \hline
        CT09 & Salvamentos de resultados & Resultados armazenados com sucesso & Funcional \\
        \hline
        CT10 & Exportação para JSON & Arquivo JSON gerado & Funcional \\
        \hline
        CT11 & Importação de mapas & Mapa carregado e reconstruído corretamente & Funcional \\
        \hline
    \end{tabular}
    \\{\small Fonte: Elaborado pelo autor.}
\end{table}

Conforme evidenciado na Tabela 2, as funcionalidades relacionadas ao gerenciamento de usuários, à criação de mundos e agentes, à visualização do histórico, ao armazenamento de resultados e às operações de importação e exportação apresentaram comportamento funcional durante os testes. Por outro lado, as funcionalidades associadas ao comportamento dos agentes, à execução das simulações e à geração de estatísticas foram classificadas como "ajustes planejados", indicando que esses componentes ainda demandam aprimoramentos no protótipo.

\subsection{Cenário Demonstrativo}

Para complementar a avaliação funcional, foi considerado um cenário demonstrativo que representa o fluxo completo de utilização da plataforma. Nesse cenário, o usuário inicia sua interação realizando o cadastro e a autenticação no sistema. Após o acesso, ele cria um mundo do Wumpus e configura um agente para atuar nesse ambiente, permitindo verificar a integração entre os recursos de gerenciamento de usuários, ambientes e agentes. Em seguida, o usuário executa uma simulação e acompanha o comportamento do agente durante sua interação com o mundo. Após a execução, são consultadas as estatísticas disponibilizadas pela plataforma para avaliar as informações produzidas durante a simulação. Os resultados obtidos são então armazenados no sistema e podem ser posteriormente recuperados por meio do histórico de execuções. Por fim, o cenário contempla a exportação dos dados do ambiente em formato JSON, seguida da importação do arquivo gerado, o que permite verificar se as informações exportadas podem ser utilizadas para reconstruir o mapa na plataforma. Dessa forma, o cenário demonstrativo reúne diferentes funcionalidades do sistema em um único fluxo, evidenciando a integração entre os módulos envolvidos.

\subsection{Discussão dos Resultados}

Os testes funcionais realizados demonstraram que a plataforma Wumpus Verse apresenta funcionamento adequado em suas principais funcionalidades de gerenciamento e persistência. Os recursos de cadastro e autenticação de usuários, criação de mundos e agentes, visualização do histórico, armazenamento de resultados e importação/exportação de cenários comportaram‑se conforme o esperado. A criação de agentes também se mostrou operacional, embora os mecanismos de personalização ainda estejam em processo de refinamento, uma vez que algumas opções de configuração necessitam de ajustes para oferecer maior flexibilidade e uma melhor experiência de utilização. Essa limitação está relacionada principalmente à necessidade de ampliar as possibilidades de configuração dos agentes disponibilizadas pela plataforma.

No que diz respeito à execução das simulações, o comportamento observado nos cenários avaliados foi considerado satisfatório, mas a tomada de decisão dos agentes e a integração com serviços externos ainda requerem aprimoramentos. Em particular, a utilização de uma API externa em sua versão gratuita pode introduzir limitações de desempenho e ocasionar eventuais inconsistências durante determinadas execuções. Outro aspecto identificado está relacionado ao módulo de estatísticas e análise dos resultados: atualmente, a plataforma disponibiliza informações básicas sobre as execuções realizadas, mas a quantidade de métricas ainda pode ser ampliada para proporcionar uma análise mais detalhada do desempenho dos agentes em diferentes cenários.

De modo geral, os resultados obtidos indicam que o protótipo atende às principais funcionalidades propostas, especialmente aquelas voltadas ao gerenciamento dos ambientes, agentes e resultados experimentais. Ao mesmo tempo, os testes permitiram identificar funcionalidades que ainda necessitam de evolução, fornecendo uma base sólida para o direcionamento dos trabalhos futuros. Entre os principais aprimoramentos planejados, destacam‑se a ampliação das opções de personalização dos agentes, o aumento da robustez da integração com serviços externos, o aperfeiçoamento dos mecanismos de tomada de decisão e a disponibilização de métricas mais detalhadas para a análise de desempenho no ambiente do Mundo do Wumpus.



